% ECONOMICS_PROBLEM: {"id": "G12-414", "title": "Equity-premium puzzle", "jel_code": "G12", "source_id": "OP-0091"}
% !TeX program = xelatex
\documentclass[11pt,a4paper]{article}
\usepackage[margin=23mm]{geometry}
\usepackage{fontspec}
\setmainfont{Latin Modern Roman}
\usepackage{amsmath,amssymb,mathtools}
\usepackage{xcolor,enumitem,booktabs,longtable,array,xurl,hyperref}
\definecolor{muted}{HTML}{53636C}
\hypersetup{colorlinks=true,urlcolor=blue,linkcolor=blue}
\setlength{\parindent}{0pt}
\setlength{\parskip}{5pt}
\newcommand{\R}{\mathbb R}
\newcommand{\E}{\mathbb E}
\newcommand{\N}{\mathbb N}
\newcommand{\ind}{\mathbf 1}
\newcommand{\Var}{\operatorname{Var}}
\newcommand{\Cov}{\operatorname{Cov}}
\newcommand{\supp}{\operatorname{supp}}
\DeclareMathOperator*{\argmin}{arg\,min}
\DeclareMathOperator*{\argmax}{arg\,max}
\newcommand{\entrymeta}[3]{{\sffamily\small\color{muted}\textbf{\texttt{#1}}\quad|\quad #2\par #3\par}}
\newcommand{\Problemref}[1]{\texttt{#1}}
\newcommand{\JELref}[2]{\texttt{#2} (JEL \texttt{#1})}
\begin{document}
\section*{Equity-premium puzzle}
\textbf{Problem ID:} \texttt{G12-414}\quad \textbf{JEL:} \texttt{G12}\par
% BEGIN ECONOMICS STATEMENT
\label{problem:OP-0091}
\label{atlas:prob:F1}

\noindent\textbf{Atlas ID:} F1; \textbf{original number:} 91; \textbf{field:} Finance. \textbf{Classification:} Editorial research formulation (theoretical/empirical); not a certified unsolved theorem.

\subsection*{Definitions and assumptions}
On a filtered probability space $(\Omega,\mathcal F,(\mathcal F_t),\mathbb P)$, let $R_{t+1}^e$ be an equity gross return and $R_t^f>0$ an $\mathcal F_t$-measurable gross risk-free return. A positive stochastic discount factor $m_{t+1}$ prices both returns if $\mathbb E_t[m_{t+1}R_{t+1}^e]=1$ and $R_t^f\mathbb E_t[m_{t+1}]=1$, where $\mathbb E_t$ denotes conditional expectation. Assume finite second moments and integrable products. Model $\theta$ specifies preferences, consumption risks, market participation and a selected equilibrium of budget-feasible household and firm decisions. For an unconstrained representative household at an interior optimum in a time-separable CRRA benchmark, $m_{t+1}=\beta(C_{t+1}/C_t)^{-\gamma}$ with $\beta\in(0,1)$, $\gamma>0$ and $C_t>0$.

\subsection*{Formal research question}
For a prespecified admissible parameter and model class $\Theta$, identify whether the joint law of consumption, portfolios, safe returns, equity returns and option prices is compatible with any $\theta\in\Theta$. The pricing restrictions imply
\[
 \mathbb E_t[R_{t+1}^e]-R_t^f
 =-\frac{\operatorname{Cov}_t(m_{t+1},R_{t+1}^e)}{\mathbb E_t[m_{t+1}]}.
\]
Quantify which restrictions reject a benchmark and which data discriminate between disaster risk, persistent consumption risk and alternative preferences. Require competing explanations to reproduce moments excluded from estimation and report the identified range of disaster probabilities or risk aversion rather than choosing parameters solely to fit the premium.

\subsection*{Interpretation and scope}
The puzzle is an empirical difficulty for a restricted benchmark, not a contradiction of no-arbitrage or of all consumption-based models. Specify return frequency, inflation adjustment, sample selection and whether the safe asset is held to maturity; a short nominal bill need not be a real risk-free asset. The observed historical premium is an estimate of a population moment, affected by sampling uncertainty and rare events. Recursive preferences require their own marginal valuation formula; the CRRA expression does not carry over unchanged. Conditional identities above do not justify replacing conditional covariances with unconditional moments without handling the time-varying safe rate. State which moments and countries are included before assessing fit, with confidence regions reflecting serial dependence.

\subsection*{Source audit and status}
[S1], also atlas link [48], defines the historical benchmark. [S2] offers disaster risk and [S3] offers long-run risk as potential resolutions. Their existence precludes calling the broad question an unsolved impossibility theorem. This is an editorial model-discrimination agenda; the audit does not certify absence of later resolutions.

\subsection*{References}

\begin{enumerate}[label={[S\arabic*]},leftmargin=*]

\item Rajnish Mehra and Edward C. Prescott (1985). \emph{The Equity Premium: A Puzzle}. Journal of Monetary Economics 15(2), 145–161. \url{https://www.sciencedirect.com/science/article/pii/0304393285900613}.
\textit{Locator and role:} Publisher or author abstract / bibliographic record; Historical US moments inconsistent with a restricted benchmark calibration; resolves link [48]. Provides background or a benchmark result; does not establish that the editorial question remains unresolved in 2026.

\item Thomas A. Rietz (1988). \emph{The Equity Risk Premium: A Solution}. Journal of Monetary Economics 22(1), 117–131. \url{https://doi.org/10.1016/0304-3932(88)90172-9}.
\textit{Locator and role:} Publisher or author abstract / bibliographic record; Low-probability consumption disasters as a potential explanation. Provides background or a benchmark result; does not establish that the editorial question remains unresolved in 2026.

\item Ravi Bansal and Amir Yaron (2004). \emph{Risks for the Long Run: A Potential Resolution of Asset Pricing Puzzles}. Journal of Finance 59(4), 1481–1509. \url{https://doi.org/10.1111/j.1540-6261.2004.00670.x}.
\textit{Locator and role:} Publisher or author abstract / bibliographic record; Long-run consumption risk and recursive-preference asset pricing. Provides background or a benchmark result; does not establish that the editorial question remains unresolved in 2026.

\end{enumerate}

\subsection*{Common conventions: Source mathematical conventions}

Unless an entry states otherwise, agent and firm sets are finite. State and action spaces are measurable, strategies and outcome maps are measurable, probability laws are specified, and expectations exist for the targets under discussion. Infinite-horizon discount factors lie in $(0,1)$ and discounted objectives are finite. Norms, tolerances and complexity input sizes must be specified for computational comparisons. Constants or primitives that appear locally are inputs, not empirical facts asserted by this catalogue.

\textbf{Identification.} A statistical model $m\in\mathcal M$ implies an observable law $P_m$ and target $\theta(m)$. At law $P$, its identified set is
\[
\Theta(P;\mathcal M)=\{\theta(m):m\in\mathcal M,\ P_m=P\}.
\]
Point identification holds when this set is a singleton, or equivalently when $P_{m_1}=P_{m_2}$ implies $\theta(m_1)=\theta(m_2)$ for all compatible models. Sharp bounds mean bounds attained, or approached when the boundary is not attained, by compatible models. Writing this set is a definition, not a solution: the substantive task is to establish informative restrictions and derive or estimate the set. An unrestricted-confounding model may give only trivial bounds.

\textbf{Inference.} Identification, estimability and valid uncertainty quantification are separate questions. Fix $\alpha\in(0,1)$. A confidence procedure $C_n$ over a declared law class $\mathcal P$ has asymptotic uniform coverage at level $1-\alpha$ when
\[
\liminf_{n\to\infty}\inf_{P\in\mathcal P}\Pr_P\{\theta(P)\in C_n\}\geq1-\alpha.
\]
For set-identified targets the coverage event must be defined separately, for example coverage of the full identified set. If an entry uses identification-indexed models instead of uniquely indexed laws, the same coverage convention is applied over those models. Pointwise limits do not establish uniform validity, and asymptotic validity does not guarantee finite-sample precision. Sample sizes, dependence structures and weak-identification sequences are part of each application.

\textbf{Counterfactuals.} $Y_i(a)$ is a potential outcome under a specified intervention, including its timing, implementation and versions. It is not $Y_i$ conditional on voluntarily choosing $a$. A full intervention vector permits interference; shorthand scalar treatment notation does not silently impose no interference. Assignment, selection, attrition, anticipation and measurement must be specified. A claim about an instrument's local effect is not automatically a population policy effect.

\textbf{Economic equilibrium.} A counterfactual model includes best responses, constraints, feasibility, market clearing, state transitions and expectations consistent with the stated information. When several equilibria exist, either a declared selector is an input or the target is set-valued. Comparisons across policy regimes must use the stated selection convention. Reduced-form relationships are not presumed invariant after an equilibrium-changing intervention.

\textbf{Optimization and welfare.} Optimization is over the stated feasible set. If attainment is not established, $\sup$ or $\inf$ and $\varepsilon$-optimal choices replace an asserted maximizer or minimizer. For example, with value $V=\sup_{a\in\mathcal A}W(a)<\infty$, an $\varepsilon$-optimal set is $\{a\in\mathcal A:W(a)\geq V-\varepsilon\}$ for $\varepsilon>0$. Benefits, costs, risk and health or environmental outcomes can be aggregated only using declared units and conversion weights. Interpersonal utility weights, the treatment of fiscal transfers and foreign profits, discounting, and the behavioral welfare criterion are explicit inputs. A comparative-static derivative additionally requires existence and appropriate differentiability of the selected counterfactual.

\textbf{Sources.} Local labels [S1], [S2], and so forth restart in every question. Locators identify the relevant abstract, model, section or result at the level actually checked; a bibliographic or abstract-level verification is not represented as inspection of an entire proof. Working papers are distinguished from later publications and changing versions. Source summaries are paraphrased. All URL checks and editorial formulations refer to the audit date stated above.
% END ECONOMICS STATEMENT
\end{document}
